\chapter{Heat-Related Illness}

\section*{Definition and Overview}
Heat-related illness is a spectrum of conditions caused by the body's failure to cope with excessive heat, ranging from mild heat exhaustion to life-threatening heat stroke. As the body heats up, it relies on sweating and increased skin blood flow to shed heat; when environmental heat and humidity overwhelm these mechanisms, core temperature rises and symptoms develop. Heat exhaustion causes fatigue, dizziness, nausea, and heavy sweating, while heat stroke, the most severe form, involves a dangerously high body temperature with confusion or collapse and requires emergency treatment.

\section*{Epidemiology}
Heat-related illness becomes more common with every heatwave, and Australia experiences some of the hottest conditions on Earth. Heatwaves cause more deaths in Australia than all other natural disasters combined. Older adults, infants, outdoor workers, athletes, people with chronic disease, and those without air-conditioning are at highest risk. Severe heat illness is fatal if untreated, and hospital presentations for heat exhaustion and heat stroke rise sharply during hot spells. Climate change is increasing both the frequency and severity of dangerous heat.

\section*{Aetiology and Pathophysiology}
\begin{itemize}
    \item \textbf{Heat exposure:} high ambient temperature, humidity, direct sun, and inadequate ventilation
    \item \textbf{Inadequate cooling:} the body's sweating and vasodilation mechanisms fail under extreme or prolonged heat
    \item \textbf{Dehydration:} loss of fluid and electrolytes through sweating reduces the body's ability to cool
    \item \textbf{Risk factors:} age extremes, chronic illness, certain medications, obesity, alcohol, and physical exertion
    \item \textbf{Heat exhaustion:} moderate heat stress with dehydration and electrolyte imbalance
    \item \textbf{Heat stroke:} core temperature above 40 degrees Celsius with failure of temperature regulation and organ dysfunction, a medical emergency
    \item \textbf{Sequelae:} high temperatures damage the brain, kidneys, liver, and other organs
\end{itemize}

\section*{Clinical Features}
\begin{itemize}
    \item \textbf{Heat exhaustion:} heavy sweating, weakness, fatigue, dizziness, headache, nausea, vomiting, muscle cramps, and rapid pulse; temperature may be normal or mildly elevated
    \item \textbf{Heat stroke:} high body temperature, hot dry or flushed skin (sweating may stop), confusion, agitation, seizures, and loss of consciousness
    \item \textbf{Heat cramps:} painful muscle spasms after exertion in heat
    \item \textbf{Heat oedema:} swelling of hands and feet in heat
    \item \textbf{Heat syncope:} fainting on standing after prolonged heat exposure
    \item \textbf{Progression:} heat exhaustion can progress to heat stroke if cooling is delayed
\end{itemize}

\section*{Differential Diagnosis}
\begin{itemize}
    \item Fever from infection can be confused with hyperthermia from heat exposure
    \item Stroke, seizure disorders, and other neurological conditions can mimic the confusion of heat stroke
    \item Hypoglycaemia and dehydration cause dizziness and weakness
    \item Drug effects, including alcohol and some medications, can cause collapse in the heat
    \item Gastroenteritis causes vomiting and collapse without heat exposure
    \item A careful history of heat exposure is essential to the diagnosis
\end{itemize}

\section*{When to Seek Medical Care}
Seek medical advice promptly for heat exhaustion that does not improve with cooling and fluids, and for anyone with risk factors who develops symptoms in the heat. Suspected heat stroke is a medical emergency.

\textbf{Call 000 (or local emergency services) for:}
\begin{itemize}
    \item Confusion, agitation, seizures, or loss of consciousness in the heat
    \item Body temperature of 40 degrees Celsius or higher
    \item Hot, dry skin with an inability to sweat
    \item Rapid breathing or pulse with collapse
    \item Failure to improve after 30 minutes of cooling and rehydration
    \item A baby, older person, or person with chronic illness who is severely unwell in the heat
\end{itemize}

\section*{Diagnosis and Investigations}
\begin{itemize}
    \item \textbf{Clinical diagnosis:} based on heat exposure and the characteristic symptoms
    \item \textbf{Temperature measurement:} core temperature is the key measurement in suspected heat stroke
    \item \textbf{Blood tests:} electrolytes, kidney function, liver function, and blood clotting in moderate to severe illness
    \item \textbf{Hydration assessment:} clinical evaluation of dehydration
    \item \textbf{Monitoring:} observation for organ dysfunction in hospital for severe cases
    \item \textbf{Risk assessment:} review of contributing factors such as medication, alcohol, and housing
\end{itemize}

\section*{Management}
\begin{itemize}
    \item \textbf{First aid:} move to shade or air-conditioning, remove excess clothing, and cool with fans, water, and cold compresses
    \item \textbf{Rehydration:} drink cool water, or oral rehydration solution where available; sports drinks for electrolyte loss
    \item \textbf{Rest:} stop activity and lie down with legs elevated
    \item \textbf{Heat cramps:} rest in a cool place, rehydrate, and gently stretch cramped muscles
    \item \textbf{Heat stroke:} immediate aggressive cooling in hospital, with intravenous fluids and organ support
    \item \textbf{Medication review:} adjust medicines that impair heat tolerance where possible
    \item \textbf{Prevention of recurrence:} advice about avoiding heat exposure and recognising early symptoms
\end{itemize}

\section*{Complications}
\begin{itemize}
    \item Heat stroke can cause permanent brain damage, kidney failure, liver injury, and rhabdomyolysis
    \item Dehydration and electrolyte disturbances, including dangerous sodium and potassium imbalances
    \item Cardiovascular stress, including heart attack in vulnerable people
    \item Death, particularly in older adults, during heatwaves
    \item Falls and injuries from fainting in the heat
    \item Worsening of chronic conditions such as diabetes, heart, lung, and kidney disease
\end{itemize}

\section*{Prognosis and Outcomes}
Heat exhaustion resolves quickly with cooling and rehydration, usually within hours, and leaves no lasting harm. The outlook for heat stroke depends on the peak temperature, its duration, and how quickly cooling begins; prompt treatment greatly improves survival, but delayed or severe cases can cause lasting organ damage or death. With sensible precautions, heat-related illness is almost entirely preventable even in extreme conditions.

\section*{Prevention and Self-Care}
\begin{itemize}
    \item Stay cool: use air-conditioning or fans, seek shade, and avoid the heat of the day
    \item Drink water regularly and don't wait for thirst
    \item Wear light, loose clothing and a hat, and use sunscreen
    \item Avoid strenuous activity during the hottest hours
    \item Never leave children, older people, or pets in parked cars
    \item Check on vulnerable neighbours, relatives, and older people during heatwaves
    \item Follow heatwave warnings and public health advice
    \item Learn first aid for heat illness so you can act quickly
\end{itemize}

\section*{Sources and Further Reading}
The following reputable sources provide further information for healthcare students and patients seeking reliable, current advice about heat-related illness.

\begin{itemize}
    \item Healthdirect. \textit{Heat-related illness: symptoms and first aid.} https://www.healthdirect.gov.au/heat-related-illness
    \item Better Health Channel. \textit{Beat the heat: preventing heat-related illness.} https://www.betterhealth.vic.gov.au/
    \item Australian Government Bureau of Meteorology. \textit{Heatwave warnings.} http://www.bom.gov.au/
    \item NHS. \textit{Heat exhaustion and heatstroke.} https://www.nhs.uk/conditions/heat-exhaustion-heatstroke/
    \item Mayo Clinic. \textit{Heatstroke: symptoms and first aid.} https://www.mayoclinic.org/
    \item World Health Organization. \textit{Heat and health.} https://www.who.int/
\end{itemize}
